% =====================================================================
%  BUỔI 19 — KẾ HOẠCH GIẢNG DẠY (từ giao-vien/01-ke-hoach-chua-de.md)
%  Sinh tự động bằng scripts/md2tex-chung.py — không sửa tay.
% =====================================================================
\documentclass[11pt, a4paper]{article}
% =====================================================================
%  CAE SHANGHAI — GÓI LỆNH DÙNG CHUNG CHO TÀI LIỆU LUYỆN THI CSCA TOÁN
%  Biên dịch bằng XeLaTeX. Mọi file .tex trong buổi học \input file này.
% =====================================================================

% ---------- NGÔN NGỮ & PHÔNG CHỮ ----------
\usepackage{fontspec}
\usepackage{polyglossia}
\setmainlanguage{vietnamese}
\setotherlanguage{english}

% Chữ Latin (Anh + Việt): TeX Gyre Termes — bản clone Times New Roman,
% đi kèm TeX Live, hỗ trợ đầy đủ dấu tiếng Việt.
\setmainfont{TeX Gyre Termes}
\setsansfont{TeX Gyre Heros}

% Chữ Hán: ưu tiên Source Han Serif SC (思源宋体) — chuẩn học thuật Trung Quốc.
% Nếu máy chưa cài, tự động dùng Songti SC (宋体, có sẵn trên macOS).
% Cả hai đều thuộc họ 宋体 nên diện mạo tài liệu không đổi.
\IfFontExistsTF{Source Han Serif SC}{
  \newcommand{\cjkfamilyname}{Source Han Serif SC}
}{
  \IfFontExistsTF{Songti SC}{
    \newcommand{\cjkfamilyname}{Songti SC}
  }{
    \newcommand{\cjkfamilyname}{Noto Serif CJK SC}
  }
}

% xeCJK tự chuyển font khi gặp ký tự Hán, không cần bọc \cjk{} thủ công.
% Đây là điều kiện bắt buộc: font Latin (TeX Gyre Termes) không có glyph chữ Hán,
% nếu không khai báo xeCJK thì mọi chữ Hán viết trần sẽ bị nuốt khỏi bản in.
\usepackage{xeCJK}
\setCJKfamilyfont{caehan}{\cjkfamilyname}[Scale=0.95]
\setCJKmainfont{\cjkfamilyname}[Scale=0.95]
\setCJKsansfont{\cjkfamilyname}[Scale=0.95]
\setCJKmonofont{\cjkfamilyname}[Scale=0.95]

% \cjk{} ép dùng font Hán trong ngữ cảnh cần thiết (ví dụ trong tiêu đề in đậm
% hoặc trong bảng có định dạng riêng). Dùng \CJKfamily chứ không phải tên font,
% vì tên font chỉ là văn bản — sẽ in ra chữ thay vì đổi font.
\newcommand{\cjk}[1]{{\CJKfamily{caehan}#1}}

% ---------- TOÁN HỌC ----------
\usepackage{amsmath, amssymb, amsthm, mathtools}
\usepackage{unicode-math}
\setmathfont{TeX Gyre Termes Math}

% ---------- BỐ CỤC TRANG ----------
% Lề phải chừa vùng an toàn của letterhead:
%   banner trên chiếm 0–19mm  → top tối thiểu 2,2cm
%   banner dưới chiếm 281–297mm → bottom tối thiểu 2,2cm
% Lề trên 2,6cm (26mm) để tiêu đề đầu trang không đè banner (banner hết ở 19mm).
\usepackage[a4paper, top=2.6cm, bottom=2.4cm, left=2.5cm, right=2.0cm]{geometry}
\usepackage{fancyhdr}
\usepackage{titlesec}
\usepackage{enumitem}
\usepackage{graphicx}
\usepackage{booktabs, array, longtable, tabularx}
\usepackage{xcolor}
\usepackage[most]{tcolorbox}
\usepackage[colorlinks=true, linkcolor=cae_blue, citecolor=cae_blue, urlcolor=cae_blue]{hyperref}

% ---------- BẢNG MÀU THƯƠNG HIỆU ----------
\definecolor{cae_blue}{RGB}{0, 51, 102}
\definecolor{cae_gold}{RGB}{196, 149, 60}
\definecolor{box_bg}{RGB}{245, 247, 250}
\definecolor{box_line}{RGB}{0, 51, 102}

% ---------- TIÊU ĐỀ MỤC ----------
% Đệm đầu trang: banner letterhead chiếm 0–19mm, mà lề trên là 22mm.
% \section* đầu trang (sau \clearpage) sát lề nên phải chừa thêm, nếu không
% sẽ đè lên banner.
\titleformat{\section}
  {\normalfont\large\bfseries\sffamily\color{cae_blue}}{\thesection.}{0.6em}{}
\titleformat{\subsection}
  {\normalfont\normalsize\bfseries\sffamily\color{cae_blue!85!black}}{\thesubsection}{0.6em}{}
\titlespacing*{\section}{0pt}{14pt}{6pt}
\titlespacing*{\subsection}{0pt}{10pt}{4pt}

% Tiêu đề không đánh số ở đầu trang: \section* sau \clearpage nằm sát lề trên
% nên đè lên banner letterhead. Dùng \vspace* ngay TRƯỚC \section* không đủ,
% vì titlesec chèn khoảng cách riêng. Cách chắc chắn là đệm bằng \null.
\newcommand{\sectiondautrang}[1]{\section*{#1}}

% Dạng bài dùng \subsection*{} nên cần định dạng riêng, đậm và có gạch chân màu
\titleformat{\subsubsection}
  {\normalfont\normalsize\bfseries\sffamily\color{cae_blue!85!black}}{\thesubsubsection}{0.6em}{}

% ---------- LETTERHEAD (ảnh nền toàn trang) ----------
% Ảnh letterhead của CAE Shanghai: banner xanh trên/dưới kèm logo chìm giữa.
% Tỉ lệ ảnh 1414x2000 = 0.7070, khớp A4 (0.7071) nên chèn thẳng không méo.
% Đặt ở lớp nền dưới cùng để chữ không bị ảnh che.
% Dùng \AddToShipoutPictureBG (không có dấu *) để áp dụng cho MỌI trang.
\usepackage{eso-pic}
\AddToShipoutPictureBG{%
  \AtPageLowerLeft{%
    \includegraphics[width=\paperwidth,height=\paperheight]{images/letterhead.png}%
  }%
}

% ---------- ĐẦU / CHÂN TRANG ----------
% Letterhead đã có banner trên và dưới, nên bỏ dòng chạy để tránh chồng chữ.
% Số trang đặt trong vùng nội dung, sát trên banner dưới, lề phải.
\pagestyle{fancy}
\setlength{\headheight}{0pt}
\fancyhf{}
\renewcommand{\headrulewidth}{0pt}
\renewcommand{\footrulewidth}{0pt}
\fancyfoot[R]{\raisebox{0.9cm}{\small\color{cae_blue}\thepage}}

% ---------- MÔI TRƯỜNG NỘI DUNG ----------
% Trong tcolorbox, LaTeX đặt parindent = parskip = 0, nên các đoạn văn trong
% khung dính liền nhau. Đặt lại khoảng cách đoạn cho mọi khung.
\newcommand{\khoangcachdoan}{\setlength{\parskip}{0.5em}\setlength{\parindent}{0pt}}
\tcbset{before upper={\khoangcachdoan}}

\newtcolorbox{dinhnghia}[1][]{
  colback=box_bg, colframe=box_line, boxrule=0.6pt, arc=1pt,
  fonttitle=\bfseries, title={Định nghĩa #1}, breakable}
\newtcolorbox{dinhly}[1][]{
  colback=box_bg, colframe=cae_gold, boxrule=0.6pt, arc=1pt,
  fonttitle=\bfseries, title={Tính chất #1}, breakable}
\newtcolorbox{vidu}[1][]{
  colback=white, colframe=black!40, boxrule=0.4pt, arc=1pt,
  fonttitle=\bfseries, title={Ví dụ #1}, breakable}
\newtcolorbox{baitap}[1][]{
  colback=white, colframe=cae_blue!50, boxrule=0.4pt, arc=1pt,
  fonttitle=\bfseries, title={Câu #1}, breakable}

% Ô đáp án: viền xanh đậm hơn, dùng trong file đáp án
\newtcolorbox{dapan}[1][]{
  colback=box_bg, colframe=cae_blue, boxrule=0.6pt, arc=1pt,
  fonttitle=\bfseries, title={Câu #1}, breakable}

% ---------- TÊN CHỨNG MINH ----------
\renewcommand{\proofname}{\normalfont\bfseries Chứng minh}

% ---------- MÔN HỌC ----------
% Đổi được ở file .tex trước khi % =====================================================================
%  CAE SHANGHAI — GÓI LỆNH DÙNG CHUNG CHO TÀI LIỆU LUYỆN THI CSCA TOÁN
%  Biên dịch bằng XeLaTeX. Mọi file .tex trong buổi học \input file này.
% =====================================================================

% ---------- NGÔN NGỮ & PHÔNG CHỮ ----------
\usepackage{fontspec}
\usepackage{polyglossia}
\setmainlanguage{vietnamese}
\setotherlanguage{english}

% Chữ Latin (Anh + Việt): TeX Gyre Termes — bản clone Times New Roman,
% đi kèm TeX Live, hỗ trợ đầy đủ dấu tiếng Việt.
\setmainfont{TeX Gyre Termes}
\setsansfont{TeX Gyre Heros}

% Chữ Hán: ưu tiên Source Han Serif SC (思源宋体) — chuẩn học thuật Trung Quốc.
% Nếu máy chưa cài, tự động dùng Songti SC (宋体, có sẵn trên macOS).
% Cả hai đều thuộc họ 宋体 nên diện mạo tài liệu không đổi.
\IfFontExistsTF{Source Han Serif SC}{
  \newcommand{\cjkfamilyname}{Source Han Serif SC}
}{
  \IfFontExistsTF{Songti SC}{
    \newcommand{\cjkfamilyname}{Songti SC}
  }{
    \newcommand{\cjkfamilyname}{Noto Serif CJK SC}
  }
}

% xeCJK tự chuyển font khi gặp ký tự Hán, không cần bọc \cjk{} thủ công.
% Đây là điều kiện bắt buộc: font Latin (TeX Gyre Termes) không có glyph chữ Hán,
% nếu không khai báo xeCJK thì mọi chữ Hán viết trần sẽ bị nuốt khỏi bản in.
\usepackage{xeCJK}
\setCJKfamilyfont{caehan}{\cjkfamilyname}[Scale=0.95]
\setCJKmainfont{\cjkfamilyname}[Scale=0.95]
\setCJKsansfont{\cjkfamilyname}[Scale=0.95]
\setCJKmonofont{\cjkfamilyname}[Scale=0.95]

% \cjk{} ép dùng font Hán trong ngữ cảnh cần thiết (ví dụ trong tiêu đề in đậm
% hoặc trong bảng có định dạng riêng). Dùng \CJKfamily chứ không phải tên font,
% vì tên font chỉ là văn bản — sẽ in ra chữ thay vì đổi font.
\newcommand{\cjk}[1]{{\CJKfamily{caehan}#1}}

% ---------- TOÁN HỌC ----------
\usepackage{amsmath, amssymb, amsthm, mathtools}
\usepackage{unicode-math}
\setmathfont{TeX Gyre Termes Math}

% ---------- BỐ CỤC TRANG ----------
% Lề phải chừa vùng an toàn của letterhead:
%   banner trên chiếm 0–19mm  → top tối thiểu 2,2cm
%   banner dưới chiếm 281–297mm → bottom tối thiểu 2,2cm
% Lề trên 2,6cm (26mm) để tiêu đề đầu trang không đè banner (banner hết ở 19mm).
\usepackage[a4paper, top=2.6cm, bottom=2.4cm, left=2.5cm, right=2.0cm]{geometry}
\usepackage{fancyhdr}
\usepackage{titlesec}
\usepackage{enumitem}
\usepackage{graphicx}
\usepackage{booktabs, array, longtable, tabularx}
\usepackage{xcolor}
\usepackage[most]{tcolorbox}
\usepackage[colorlinks=true, linkcolor=cae_blue, citecolor=cae_blue, urlcolor=cae_blue]{hyperref}

% ---------- BẢNG MÀU THƯƠNG HIỆU ----------
\definecolor{cae_blue}{RGB}{0, 51, 102}
\definecolor{cae_gold}{RGB}{196, 149, 60}
\definecolor{box_bg}{RGB}{245, 247, 250}
\definecolor{box_line}{RGB}{0, 51, 102}

% ---------- TIÊU ĐỀ MỤC ----------
% Đệm đầu trang: banner letterhead chiếm 0–19mm, mà lề trên là 22mm.
% \section* đầu trang (sau \clearpage) sát lề nên phải chừa thêm, nếu không
% sẽ đè lên banner.
\titleformat{\section}
  {\normalfont\large\bfseries\sffamily\color{cae_blue}}{\thesection.}{0.6em}{}
\titleformat{\subsection}
  {\normalfont\normalsize\bfseries\sffamily\color{cae_blue!85!black}}{\thesubsection}{0.6em}{}
\titlespacing*{\section}{0pt}{14pt}{6pt}
\titlespacing*{\subsection}{0pt}{10pt}{4pt}

% Tiêu đề không đánh số ở đầu trang: \section* sau \clearpage nằm sát lề trên
% nên đè lên banner letterhead. Dùng \vspace* ngay TRƯỚC \section* không đủ,
% vì titlesec chèn khoảng cách riêng. Cách chắc chắn là đệm bằng \null.
\newcommand{\sectiondautrang}[1]{\section*{#1}}

% Dạng bài dùng \subsection*{} nên cần định dạng riêng, đậm và có gạch chân màu
\titleformat{\subsubsection}
  {\normalfont\normalsize\bfseries\sffamily\color{cae_blue!85!black}}{\thesubsubsection}{0.6em}{}

% ---------- LETTERHEAD (ảnh nền toàn trang) ----------
% Ảnh letterhead của CAE Shanghai: banner xanh trên/dưới kèm logo chìm giữa.
% Tỉ lệ ảnh 1414x2000 = 0.7070, khớp A4 (0.7071) nên chèn thẳng không méo.
% Đặt ở lớp nền dưới cùng để chữ không bị ảnh che.
% Dùng \AddToShipoutPictureBG (không có dấu *) để áp dụng cho MỌI trang.
\usepackage{eso-pic}
\AddToShipoutPictureBG{%
  \AtPageLowerLeft{%
    \includegraphics[width=\paperwidth,height=\paperheight]{images/letterhead.png}%
  }%
}

% ---------- ĐẦU / CHÂN TRANG ----------
% Letterhead đã có banner trên và dưới, nên bỏ dòng chạy để tránh chồng chữ.
% Số trang đặt trong vùng nội dung, sát trên banner dưới, lề phải.
\pagestyle{fancy}
\setlength{\headheight}{0pt}
\fancyhf{}
\renewcommand{\headrulewidth}{0pt}
\renewcommand{\footrulewidth}{0pt}
\fancyfoot[R]{\raisebox{0.9cm}{\small\color{cae_blue}\thepage}}

% ---------- MÔI TRƯỜNG NỘI DUNG ----------
% Trong tcolorbox, LaTeX đặt parindent = parskip = 0, nên các đoạn văn trong
% khung dính liền nhau. Đặt lại khoảng cách đoạn cho mọi khung.
\newcommand{\khoangcachdoan}{\setlength{\parskip}{0.5em}\setlength{\parindent}{0pt}}
\tcbset{before upper={\khoangcachdoan}}

\newtcolorbox{dinhnghia}[1][]{
  colback=box_bg, colframe=box_line, boxrule=0.6pt, arc=1pt,
  fonttitle=\bfseries, title={Định nghĩa #1}, breakable}
\newtcolorbox{dinhly}[1][]{
  colback=box_bg, colframe=cae_gold, boxrule=0.6pt, arc=1pt,
  fonttitle=\bfseries, title={Tính chất #1}, breakable}
\newtcolorbox{vidu}[1][]{
  colback=white, colframe=black!40, boxrule=0.4pt, arc=1pt,
  fonttitle=\bfseries, title={Ví dụ #1}, breakable}
\newtcolorbox{baitap}[1][]{
  colback=white, colframe=cae_blue!50, boxrule=0.4pt, arc=1pt,
  fonttitle=\bfseries, title={Câu #1}, breakable}

% Ô đáp án: viền xanh đậm hơn, dùng trong file đáp án
\newtcolorbox{dapan}[1][]{
  colback=box_bg, colframe=cae_blue, boxrule=0.6pt, arc=1pt,
  fonttitle=\bfseries, title={Câu #1}, breakable}

% ---------- TÊN CHỨNG MINH ----------
\renewcommand{\proofname}{\normalfont\bfseries Chứng minh}

% ---------- MÔN HỌC ----------
% Đổi được ở file .tex trước khi % =====================================================================
%  CAE SHANGHAI — GÓI LỆNH DÙNG CHUNG CHO TÀI LIỆU LUYỆN THI CSCA TOÁN
%  Biên dịch bằng XeLaTeX. Mọi file .tex trong buổi học \input file này.
% =====================================================================

% ---------- NGÔN NGỮ & PHÔNG CHỮ ----------
\usepackage{fontspec}
\usepackage{polyglossia}
\setmainlanguage{vietnamese}
\setotherlanguage{english}

% Chữ Latin (Anh + Việt): TeX Gyre Termes — bản clone Times New Roman,
% đi kèm TeX Live, hỗ trợ đầy đủ dấu tiếng Việt.
\setmainfont{TeX Gyre Termes}
\setsansfont{TeX Gyre Heros}

% Chữ Hán: ưu tiên Source Han Serif SC (思源宋体) — chuẩn học thuật Trung Quốc.
% Nếu máy chưa cài, tự động dùng Songti SC (宋体, có sẵn trên macOS).
% Cả hai đều thuộc họ 宋体 nên diện mạo tài liệu không đổi.
\IfFontExistsTF{Source Han Serif SC}{
  \newcommand{\cjkfamilyname}{Source Han Serif SC}
}{
  \IfFontExistsTF{Songti SC}{
    \newcommand{\cjkfamilyname}{Songti SC}
  }{
    \newcommand{\cjkfamilyname}{Noto Serif CJK SC}
  }
}

% xeCJK tự chuyển font khi gặp ký tự Hán, không cần bọc \cjk{} thủ công.
% Đây là điều kiện bắt buộc: font Latin (TeX Gyre Termes) không có glyph chữ Hán,
% nếu không khai báo xeCJK thì mọi chữ Hán viết trần sẽ bị nuốt khỏi bản in.
\usepackage{xeCJK}
\setCJKfamilyfont{caehan}{\cjkfamilyname}[Scale=0.95]
\setCJKmainfont{\cjkfamilyname}[Scale=0.95]
\setCJKsansfont{\cjkfamilyname}[Scale=0.95]
\setCJKmonofont{\cjkfamilyname}[Scale=0.95]

% \cjk{} ép dùng font Hán trong ngữ cảnh cần thiết (ví dụ trong tiêu đề in đậm
% hoặc trong bảng có định dạng riêng). Dùng \CJKfamily chứ không phải tên font,
% vì tên font chỉ là văn bản — sẽ in ra chữ thay vì đổi font.
\newcommand{\cjk}[1]{{\CJKfamily{caehan}#1}}

% ---------- TOÁN HỌC ----------
\usepackage{amsmath, amssymb, amsthm, mathtools}
\usepackage{unicode-math}
\setmathfont{TeX Gyre Termes Math}

% ---------- BỐ CỤC TRANG ----------
% Lề phải chừa vùng an toàn của letterhead:
%   banner trên chiếm 0–19mm  → top tối thiểu 2,2cm
%   banner dưới chiếm 281–297mm → bottom tối thiểu 2,2cm
% Lề trên 2,6cm (26mm) để tiêu đề đầu trang không đè banner (banner hết ở 19mm).
\usepackage[a4paper, top=2.6cm, bottom=2.4cm, left=2.5cm, right=2.0cm]{geometry}
\usepackage{fancyhdr}
\usepackage{titlesec}
\usepackage{enumitem}
\usepackage{graphicx}
\usepackage{booktabs, array, longtable, tabularx}
\usepackage{xcolor}
\usepackage[most]{tcolorbox}
\usepackage[colorlinks=true, linkcolor=cae_blue, citecolor=cae_blue, urlcolor=cae_blue]{hyperref}

% ---------- BẢNG MÀU THƯƠNG HIỆU ----------
\definecolor{cae_blue}{RGB}{0, 51, 102}
\definecolor{cae_gold}{RGB}{196, 149, 60}
\definecolor{box_bg}{RGB}{245, 247, 250}
\definecolor{box_line}{RGB}{0, 51, 102}

% ---------- TIÊU ĐỀ MỤC ----------
% Đệm đầu trang: banner letterhead chiếm 0–19mm, mà lề trên là 22mm.
% \section* đầu trang (sau \clearpage) sát lề nên phải chừa thêm, nếu không
% sẽ đè lên banner.
\titleformat{\section}
  {\normalfont\large\bfseries\sffamily\color{cae_blue}}{\thesection.}{0.6em}{}
\titleformat{\subsection}
  {\normalfont\normalsize\bfseries\sffamily\color{cae_blue!85!black}}{\thesubsection}{0.6em}{}
\titlespacing*{\section}{0pt}{14pt}{6pt}
\titlespacing*{\subsection}{0pt}{10pt}{4pt}

% Tiêu đề không đánh số ở đầu trang: \section* sau \clearpage nằm sát lề trên
% nên đè lên banner letterhead. Dùng \vspace* ngay TRƯỚC \section* không đủ,
% vì titlesec chèn khoảng cách riêng. Cách chắc chắn là đệm bằng \null.
\newcommand{\sectiondautrang}[1]{\section*{#1}}

% Dạng bài dùng \subsection*{} nên cần định dạng riêng, đậm và có gạch chân màu
\titleformat{\subsubsection}
  {\normalfont\normalsize\bfseries\sffamily\color{cae_blue!85!black}}{\thesubsubsection}{0.6em}{}

% ---------- LETTERHEAD (ảnh nền toàn trang) ----------
% Ảnh letterhead của CAE Shanghai: banner xanh trên/dưới kèm logo chìm giữa.
% Tỉ lệ ảnh 1414x2000 = 0.7070, khớp A4 (0.7071) nên chèn thẳng không méo.
% Đặt ở lớp nền dưới cùng để chữ không bị ảnh che.
% Dùng \AddToShipoutPictureBG (không có dấu *) để áp dụng cho MỌI trang.
\usepackage{eso-pic}
\AddToShipoutPictureBG{%
  \AtPageLowerLeft{%
    \includegraphics[width=\paperwidth,height=\paperheight]{images/letterhead.png}%
  }%
}

% ---------- ĐẦU / CHÂN TRANG ----------
% Letterhead đã có banner trên và dưới, nên bỏ dòng chạy để tránh chồng chữ.
% Số trang đặt trong vùng nội dung, sát trên banner dưới, lề phải.
\pagestyle{fancy}
\setlength{\headheight}{0pt}
\fancyhf{}
\renewcommand{\headrulewidth}{0pt}
\renewcommand{\footrulewidth}{0pt}
\fancyfoot[R]{\raisebox{0.9cm}{\small\color{cae_blue}\thepage}}

% ---------- MÔI TRƯỜNG NỘI DUNG ----------
% Trong tcolorbox, LaTeX đặt parindent = parskip = 0, nên các đoạn văn trong
% khung dính liền nhau. Đặt lại khoảng cách đoạn cho mọi khung.
\newcommand{\khoangcachdoan}{\setlength{\parskip}{0.5em}\setlength{\parindent}{0pt}}
\tcbset{before upper={\khoangcachdoan}}

\newtcolorbox{dinhnghia}[1][]{
  colback=box_bg, colframe=box_line, boxrule=0.6pt, arc=1pt,
  fonttitle=\bfseries, title={Định nghĩa #1}, breakable}
\newtcolorbox{dinhly}[1][]{
  colback=box_bg, colframe=cae_gold, boxrule=0.6pt, arc=1pt,
  fonttitle=\bfseries, title={Tính chất #1}, breakable}
\newtcolorbox{vidu}[1][]{
  colback=white, colframe=black!40, boxrule=0.4pt, arc=1pt,
  fonttitle=\bfseries, title={Ví dụ #1}, breakable}
\newtcolorbox{baitap}[1][]{
  colback=white, colframe=cae_blue!50, boxrule=0.4pt, arc=1pt,
  fonttitle=\bfseries, title={Câu #1}, breakable}

% Ô đáp án: viền xanh đậm hơn, dùng trong file đáp án
\newtcolorbox{dapan}[1][]{
  colback=box_bg, colframe=cae_blue, boxrule=0.6pt, arc=1pt,
  fonttitle=\bfseries, title={Câu #1}, breakable}

% ---------- TÊN CHỨNG MINH ----------
\renewcommand{\proofname}{\normalfont\bfseries Chứng minh}

% ---------- MÔN HỌC ----------
% Đổi được ở file .tex trước khi \input{cae-style} bằng \renewcommand{\monhoc}{...}
\providecommand{\monhoc}{Môn Toán}

% ---------- TRANG BÌA GỌN ----------
% #1 = loại tài liệu (TÀI LIỆU CHUẨN BỊ, TÀI LIỆU BUỔI HỌC, ...)
% #2 = số buổi
% #3 = tên buổi tiếng Việt
% #4 = tên buổi tiếng Hán
% #5 = module và tỉ trọng
\newcommand{\trangbia}[5]{%
  \begin{center}
    \vspace*{0.6cm}
    {\large\bfseries\color{cae_blue} CAE SHANGHAI}\\[0.15cm]
    {\footnotesize Tài liệu luyện thi CSCA · \monhoc}\\[0.5cm]
    {\color{cae_gold}\rule{0.55\linewidth}{0.8pt}}\\[0.35cm]
    {\small\bfseries\color{cae_blue!80!black} #1}\\[0.35cm]
    {\Large\bfseries\color{cae_blue} BUỔI #2 — #3}\\[0.2cm]
    {\large\color{cae_blue!85!black} \cjk{#4}}\\[0.25cm]
    \ifx\relax#5\relax\else{\small #5}\\[0.35cm]\fi
    {\color{cae_gold}\rule{0.55\linewidth}{0.8pt}}\\[0.3cm]
    {\footnotesize Biên soạn: Hoàng Dương, CAE Shanghai \quad·\quad \today}
  \end{center}
  \vspace{0.3cm}
}

% ---------- TRANG BÌA CHO ĐỀ KIỂM TRA ----------
% Khác \trangbia ở chỗ không có chữ "BUỔI N" — dùng cho đề thi, không phải bài giảng.
% #1 = loại tài liệu, #2 = tên kỳ thi (Việt), #3 = tên kỳ thi (Hán), #4 = thông số
\newcommand{\biade}[4]{%
  \begin{center}
    \vspace*{0.6cm}
    {\large\bfseries\color{cae_blue} CAE SHANGHAI}\\[0.15cm]
    {\footnotesize Tài liệu luyện thi CSCA · \monhoc}\\[0.5cm]
    {\color{cae_gold}\rule{0.55\linewidth}{0.8pt}}\\[0.35cm]
    {\small\bfseries\color{cae_blue!80!black} #1}\\[0.35cm]
    {\Large\bfseries\color{cae_blue} #2}\\[0.2cm]
    {\large\color{cae_blue!85!black} \cjk{#3}}\\[0.25cm]
    {\small #4}\\[0.35cm]
    {\color{cae_gold}\rule{0.55\linewidth}{0.8pt}}\\[0.3cm]
    {\footnotesize Biên soạn: Hoàng Dương, CAE Shanghai \quad·\quad \today}
  \end{center}
  \vspace{0.3cm}
}

% ---------- TIỆN ÍCH ----------
% Nhãn từ khóa in đậm ở đầu đoạn, dùng cho Từ vựng / Từ khóa nhận dạng / Chú ý
% Khoảng cách đoạn do khung tự đặt (xem \khoangcachdoan ở trên),
% ở đây chỉ cần tách đoạn và bỏ thụt đầu dòng.
\newcommand{\nhan}[1]{\par\noindent\textbf{#1.}\ \ignorespaces}

% Bảng gọn, không đánh số
\newcommand{\bang}[1]{%
  \begin{center}#1\end{center}}

% Đánh dấu ô trống cho học sinh điền
\newcommand{\chotrong}{\makebox[3.2cm]{\dotfill}}
 bằng \renewcommand{\monhoc}{...}
\providecommand{\monhoc}{Môn Toán}

% ---------- TRANG BÌA GỌN ----------
% #1 = loại tài liệu (TÀI LIỆU CHUẨN BỊ, TÀI LIỆU BUỔI HỌC, ...)
% #2 = số buổi
% #3 = tên buổi tiếng Việt
% #4 = tên buổi tiếng Hán
% #5 = module và tỉ trọng
\newcommand{\trangbia}[5]{%
  \begin{center}
    \vspace*{0.6cm}
    {\large\bfseries\color{cae_blue} CAE SHANGHAI}\\[0.15cm]
    {\footnotesize Tài liệu luyện thi CSCA · \monhoc}\\[0.5cm]
    {\color{cae_gold}\rule{0.55\linewidth}{0.8pt}}\\[0.35cm]
    {\small\bfseries\color{cae_blue!80!black} #1}\\[0.35cm]
    {\Large\bfseries\color{cae_blue} BUỔI #2 — #3}\\[0.2cm]
    {\large\color{cae_blue!85!black} \cjk{#4}}\\[0.25cm]
    \ifx\relax#5\relax\else{\small #5}\\[0.35cm]\fi
    {\color{cae_gold}\rule{0.55\linewidth}{0.8pt}}\\[0.3cm]
    {\footnotesize Biên soạn: Hoàng Dương, CAE Shanghai \quad·\quad \today}
  \end{center}
  \vspace{0.3cm}
}

% ---------- TRANG BÌA CHO ĐỀ KIỂM TRA ----------
% Khác \trangbia ở chỗ không có chữ "BUỔI N" — dùng cho đề thi, không phải bài giảng.
% #1 = loại tài liệu, #2 = tên kỳ thi (Việt), #3 = tên kỳ thi (Hán), #4 = thông số
\newcommand{\biade}[4]{%
  \begin{center}
    \vspace*{0.6cm}
    {\large\bfseries\color{cae_blue} CAE SHANGHAI}\\[0.15cm]
    {\footnotesize Tài liệu luyện thi CSCA · \monhoc}\\[0.5cm]
    {\color{cae_gold}\rule{0.55\linewidth}{0.8pt}}\\[0.35cm]
    {\small\bfseries\color{cae_blue!80!black} #1}\\[0.35cm]
    {\Large\bfseries\color{cae_blue} #2}\\[0.2cm]
    {\large\color{cae_blue!85!black} \cjk{#3}}\\[0.25cm]
    {\small #4}\\[0.35cm]
    {\color{cae_gold}\rule{0.55\linewidth}{0.8pt}}\\[0.3cm]
    {\footnotesize Biên soạn: Hoàng Dương, CAE Shanghai \quad·\quad \today}
  \end{center}
  \vspace{0.3cm}
}

% ---------- TIỆN ÍCH ----------
% Nhãn từ khóa in đậm ở đầu đoạn, dùng cho Từ vựng / Từ khóa nhận dạng / Chú ý
% Khoảng cách đoạn do khung tự đặt (xem \khoangcachdoan ở trên),
% ở đây chỉ cần tách đoạn và bỏ thụt đầu dòng.
\newcommand{\nhan}[1]{\par\noindent\textbf{#1.}\ \ignorespaces}

% Bảng gọn, không đánh số
\newcommand{\bang}[1]{%
  \begin{center}#1\end{center}}

% Đánh dấu ô trống cho học sinh điền
\newcommand{\chotrong}{\makebox[3.2cm]{\dotfill}}
 bằng \renewcommand{\monhoc}{...}
\providecommand{\monhoc}{Môn Toán}

% ---------- TRANG BÌA GỌN ----------
% #1 = loại tài liệu (TÀI LIỆU CHUẨN BỊ, TÀI LIỆU BUỔI HỌC, ...)
% #2 = số buổi
% #3 = tên buổi tiếng Việt
% #4 = tên buổi tiếng Hán
% #5 = module và tỉ trọng
\newcommand{\trangbia}[5]{%
  \begin{center}
    \vspace*{0.6cm}
    {\large\bfseries\color{cae_blue} CAE SHANGHAI}\\[0.15cm]
    {\footnotesize Tài liệu luyện thi CSCA · \monhoc}\\[0.5cm]
    {\color{cae_gold}\rule{0.55\linewidth}{0.8pt}}\\[0.35cm]
    {\small\bfseries\color{cae_blue!80!black} #1}\\[0.35cm]
    {\Large\bfseries\color{cae_blue} BUỔI #2 — #3}\\[0.2cm]
    {\large\color{cae_blue!85!black} \cjk{#4}}\\[0.25cm]
    \ifx\relax#5\relax\else{\small #5}\\[0.35cm]\fi
    {\color{cae_gold}\rule{0.55\linewidth}{0.8pt}}\\[0.3cm]
    {\footnotesize Biên soạn: Hoàng Dương, CAE Shanghai \quad·\quad \today}
  \end{center}
  \vspace{0.3cm}
}

% ---------- TRANG BÌA CHO ĐỀ KIỂM TRA ----------
% Khác \trangbia ở chỗ không có chữ "BUỔI N" — dùng cho đề thi, không phải bài giảng.
% #1 = loại tài liệu, #2 = tên kỳ thi (Việt), #3 = tên kỳ thi (Hán), #4 = thông số
\newcommand{\biade}[4]{%
  \begin{center}
    \vspace*{0.6cm}
    {\large\bfseries\color{cae_blue} CAE SHANGHAI}\\[0.15cm]
    {\footnotesize Tài liệu luyện thi CSCA · \monhoc}\\[0.5cm]
    {\color{cae_gold}\rule{0.55\linewidth}{0.8pt}}\\[0.35cm]
    {\small\bfseries\color{cae_blue!80!black} #1}\\[0.35cm]
    {\Large\bfseries\color{cae_blue} #2}\\[0.2cm]
    {\large\color{cae_blue!85!black} \cjk{#3}}\\[0.25cm]
    {\small #4}\\[0.35cm]
    {\color{cae_gold}\rule{0.55\linewidth}{0.8pt}}\\[0.3cm]
    {\footnotesize Biên soạn: Hoàng Dương, CAE Shanghai \quad·\quad \today}
  \end{center}
  \vspace{0.3cm}
}

% ---------- TIỆN ÍCH ----------
% Nhãn từ khóa in đậm ở đầu đoạn, dùng cho Từ vựng / Từ khóa nhận dạng / Chú ý
% Khoảng cách đoạn do khung tự đặt (xem \khoangcachdoan ở trên),
% ở đây chỉ cần tách đoạn và bỏ thụt đầu dòng.
\newcommand{\nhan}[1]{\par\noindent\textbf{#1.}\ \ignorespaces}

% Bảng gọn, không đánh số
\newcommand{\bang}[1]{%
  \begin{center}#1\end{center}}

% Đánh dấu ô trống cho học sinh điền
\newcommand{\chotrong}{\makebox[3.2cm]{\dotfill}}


\begin{document}

\trangbia{KẾ HOẠCH GIẢNG DẠY}{19}{Toàn bộ chương trình}{综合模拟}{}

\section{PHÂN BỔ THỜI GIAN}

\bang{
\begin{tabularx}{\linewidth}{@{\hspace{6pt}}Xcc@{\hspace{6pt}}}
\toprule
\textbf{Hoạt động} & \textbf{Tỉ trọng} & \textbf{Thời lượng} \\
\midrule
Tổng quan và phân loại lỗi & 10\% & 8 phút \\
Chữa nhanh các câu đúng & 25\% & 20 phút \\
Chữa kỹ các câu sai nhiều & 45\% & 36 phút \\
Hệ thống lại dạng bài trọng tâm và kỹ thuật xử lý nhanh & 20\% & 16 phút \\
\bottomrule
\end{tabularx}}

Đây là buổi chữa đề tổng hợp thứ hai, phủ toàn bộ chương trình. Khác với các buổi chữa đề theo chuyên đề, buổi này không dạy kiến thức mới mà \textbf{chẩn đoán lỗi hệ thống}: học sinh đã học hết 20 buổi, vấn đề còn lại là tốc độ và độ chính xác khi trộn lẫn mọi dạng bài.

\section{BẢNG PHÂN BỐ CÂU THEO NHÓM NỘI DUNG}

\bang{
\begin{tabularx}{\linewidth}{@{\hspace{6pt}}Xlcc@{\hspace{6pt}}}
\toprule
\textbf{Nhóm nội dung} & \textbf{Câu} & \textbf{Số câu} & \textbf{Tỉ lệ} \\
\midrule
Lượng giác & 1 đến 10 & 10 & 20.8\% \\
Hình học giải tích & 11 đến 19 & 9 & 18.8\% \\
Dãy số & 20 đến 25 & 6 & 12.5\% \\
Hàm số & 26 đến 31 & 6 & 12.5\% \\
Tập hợp và bất đẳng thức & 32 đến 36 & 5 & 10.4\% \\
Mũ và logarit & 37 đến 40 & 4 & 8.3\% \\
Vector và số phức & 41 đến 43 & 3 & 6.3\% \\
Xác suất và thống kê & 44 đến 45 & 2 & 4.2\% \\
Hình học phẳng và không gian & 46 đến 48 & 3 & 6.3\% \\
\bottomrule
\end{tabularx}}

Phân bố này khớp với trọng số đề thi thật. Lượng giác và Hình học giải tích chiếm gần 40\% số câu, nên hai nhóm này quyết định điểm số.

\section{BẢNG ĐÁP ÁN}

\bang{
\begin{tabular}{@{}lcccccccccccc@{}}
\toprule
\textbf{Câu} & \textbf{1} & \textbf{2} & \textbf{3} & \textbf{4} & \textbf{5} & \textbf{6} & \textbf{7} & \textbf{8} & \textbf{9} & \textbf{10} & \textbf{11} & \textbf{12} \\
\midrule
Đáp án & B & A & A & B & B & A & B & B & D & D & A & D \\
\bottomrule
\end{tabular}}

\bang{
\begin{tabular}{@{}lcccccccccccc@{}}
\toprule
\textbf{Câu} & \textbf{13} & \textbf{14} & \textbf{15} & \textbf{16} & \textbf{17} & \textbf{18} & \textbf{19} & \textbf{20} & \textbf{21} & \textbf{22} & \textbf{23} & \textbf{24} \\
\midrule
Đáp án & B & D & B & C & C & D & D & C & C & D & C & A \\
\bottomrule
\end{tabular}}

\bang{
\begin{tabular}{@{}lcccccccccccc@{}}
\toprule
\textbf{Câu} & \textbf{25} & \textbf{26} & \textbf{27} & \textbf{28} & \textbf{29} & \textbf{30} & \textbf{31} & \textbf{32} & \textbf{33} & \textbf{34} & \textbf{35} & \textbf{36} \\
\midrule
Đáp án & A & C & A & D & D & A & C & A & C & D & C & A \\
\bottomrule
\end{tabular}}

\bang{
\begin{tabular}{@{}lcccccccccccc@{}}
\toprule
\textbf{Câu} & \textbf{37} & \textbf{38} & \textbf{39} & \textbf{40} & \textbf{41} & \textbf{42} & \textbf{43} & \textbf{44} & \textbf{45} & \textbf{46} & \textbf{47} & \textbf{48} \\
\midrule
Đáp án & B & B & C & B & C & B & A & A & D & B & C & D \\
\bottomrule
\end{tabular}}

\section{KHỐI 1 — TỔNG QUAN VÀ PHÂN LOẠI LỖI (8 phút)}

\subsection{Mục tiêu}

Học sinh thấy được bức tranh tổng thể: mình mất điểm ở nhóm nào, vì lỗi gì, và lỗi đó có sửa được trong thời gian còn lại hay không.

\subsection{Kịch bản}

\textbf{Phút 1–3 · Công bố phổ điểm.} Thu bài, thống kê nhanh ba con số lên bảng: điểm trung bình, điểm cao nhất, điểm thấp nhất. Không đọc tên từng người.

\textbf{Phút 4–6 · Phân loại lỗi thành ba nhóm.} Ghi ba tiêu đề lên bảng và xếp từng câu sai vào một trong ba:

\bang{
\begin{tabularx}{\linewidth}{@{\hspace{6pt}}lXl@{\hspace{6pt}}}
\toprule
\textbf{Loại lỗi} & \textbf{Dấu hiệu} & \textbf{Xử lý} \\
\midrule
Lỗi kiến thức & Không biết công thức, chưa từng gặp dạng & Phải học lại, cần thời gian \\
Lỗi kỹ thuật & Biết công thức nhưng biến đổi sai dấu, sai hệ số & Sửa được ngay trong buổi này \\
Lỗi đọc đề & Nhầm từ khóa, hiểu sai yêu cầu, hết thời gian & Sửa được, quan trọng nhất \\
\bottomrule
\end{tabularx}}

\textbf{Phút 7–8 · Chốt thông điệp.} Phần lớn điểm bị mất ở hai loại sau, không phải loại đầu. Nói rõ: trong 72 phút còn lại, lớp sẽ chữa kỹ đúng những câu thuộc hai loại sau.

\subsection{Điểm cần chốt}

Học sinh phải tự đếm được số câu mình sai do lỗi kỹ thuật và lỗi đọc đề. Đây là con số các em cải thiện được ngay trong kỳ thi cuối kỳ.

\section{KHỐI 2 — CHỮA NHANH CÁC CÂU ĐÚNG (20 phút)}

\subsection{Mục tiêu}

Củng cố phản xạ đúng, đồng thời rút ra \textbf{đường đi ngắn nhất} cho mỗi dạng. Học sinh đã làm đúng không cần nghe lại lời giải, chỉ cần nghe cách làm nhanh hơn.

\subsection{Kịch bản}

Chữa theo cụm, mỗi cụm 3 đến 4 câu cùng dạng, khoảng 2 phút mỗi cụm. Với mỗi câu, chỉ trả lời một câu hỏi: \textbf{"Câu này làm trong bao nhiêu giây, bằng cách nào?"}

\bang{
\begin{tabularx}{\linewidth}{@{\hspace{6pt}}llX@{\hspace{6pt}}}
\toprule
\textbf{Cụm} & \textbf{Câu} & \textbf{Trọng tâm} \\
\midrule
Giá trị lượng giác cơ bản & 1, 2, 7 & Dấu theo góc phần tư, công thức góc bù \\
Công thức nhân đôi & 8, 9, 10 & $\sin 2\alpha = 2\sin\alpha\cos\alpha$, $1 - \cos 2\alpha = 2\sin^2\alpha$ \\
Biến đổi theo $\tan$ & 5 & Chia tử và mẫu cho $\cos$ để quy về $\tan$ \\
Dạng $a\sin x + b\cos x$ & 6 & Nhớ ngay $\sqrt{a^2+b^2}$, không biến đổi \\
Chu kì và giá trị lớn nhất & 3, 4 & Công thức $T = \dfrac{2\pi}{\lvert \omega \rvert}$, miền giá trị của $\sin$ \\
Đường thẳng & 11, 12, 19 & Đọc hệ số góc, điều kiện song song và vuông góc \\
Đường tròn & 13, 18 & Tâm $\left(-\dfrac{D}{2}, -\dfrac{E}{2}\right)$, công thức bán kính \\
Khoảng cách & 14 & Thay số vào công thức khoảng cách điểm đến đường thẳng \\
Conic & 15, 16, 17 & Đọc thẳng $p$, $a$, $b$, $c$ từ phương trình chuẩn \\
Cấp số cộng và nhân & 20, 21, 22, 23, 24 & Công thức số hạng tổng quát và tổng $n$ số hạng đầu \\
Hệ thức truy hồi & 25 & Tính lần lượt, không tìm công thức tổng quát \\
Tập xác định và cực trị & 26, 27, 30 & Điều kiện căn và mẫu, trục đối xứng parabol \\
Đạo hàm và giá trị lớn nhất & 28, 31 & Xét dấu đạo hàm, phương trình có nghiệm \\
Mũ và logarit & 37, 38, 39, 40 & Cùng cơ số, đổi chiều khi cơ số nhỏ hơn $1$ \\
Vector và số phức & 41, 42, 43 & Tích vô hướng, cộng số phức, môđun \\
Xác suất và thống kê & 44, 45 & Đếm kết quả thuận lợi, trung vị \\
Tập hợp và bất đẳng thức & 32, 33, 34, 35, 36 & Giao tập hợp, giá trị tuyệt đối, Cauchy \\
Hình học không gian & 46, 47, 48 & Tâm mặt cầu, khoảng cách, diện tích mặt cầu \\
\bottomrule
\end{tabularx}}

\subsection{Lỗi học sinh hay mắc — cần cảnh báo trước}

\bang{
\begin{tabularx}{\linewidth}{@{\hspace{6pt}}lX@{\hspace{6pt}}}
\toprule
\textbf{Lỗi} & \textbf{Cách cảnh báo} \\
\midrule
Quên dấu khi khai căn theo góc phần tư & Nói trước: "Góc phần tư thứ ba thì $\cos$ âm, không được ghi $\pm$" \\
Nhầm chiều bất đẳng thức mũ và logarit & Viết lên bảng: cơ số lớn hơn $1$ giữ chiều, nhỏ hơn $1$ đổi chiều \\
Lấy $\Delta \ge 0$ thay vì $\Delta \le 0$ ở bài cực trị & Nhắc: "Điều kiện có nghiệm là $\Delta \ge 0$, không phải ngược lại" \\
Nhầm trung vị với trung bình cộng & Nhấn: "Trung vị là giá trị đứng giữa sau khi sắp xếp, không phải tổng chia số lượng" \\
Nhầm diện tích mặt cầu với thể tích & Viết lên bảng: $S = 4\pi R^2$, $V = \dfrac{4}{3}\pi R^3$ \\
\bottomrule
\end{tabularx}}

\section{KHỐI 3 — CHỮA KỸ CÁC CÂU SAI NHIỀU (36 phút)}

\subsection{Mục tiêu}

Mỗi câu sai nhiều phải được chữa đến mức học sinh \textbf{tự làm lại được một câu tương tự} ngay tại lớp, không chỉ hiểu lời giải.

\subsection{Thứ tự ưu tiên}

\bang{
\begin{tabularx}{\linewidth}{@{\hspace{6pt}}clX@{\hspace{6pt}}}
\toprule
\textbf{Ưu tiên} & \textbf{Câu} & \textbf{Lý do} \\
\midrule
1 & 31 & Tìm giá trị lớn nhất của hàm phân thức, phải đưa về điều kiện có nghiệm \\
2 & 26, 34 & Điều kiện xác định và giá trị tuyệt đối, học sinh hay quên loại nghiệm \\
3 & 40 & Bất đẳng thức mũ có cơ số nhỏ hơn $1$, phải đổi chiều \\
4 & 19 & Điều kiện vuông góc rồi viết phương trình, hai bước dễ sai dấu \\
5 & 24, 25 & Cấp số nhân cho gián tiếp và hệ thức truy hồi \\
6 & 28, 30 & Đạo hàm xét dấu và điều kiện logarit \\
7 & 8 & Bình phương đẳng thức $\sin\alpha + \cos\alpha$ để ra $\sin 2\alpha$ \\
8 & 9, 10 & Công thức nhân đôi và hạ bậc \\
9 & 36 & Bất đẳng thức Cauchy, hay quên điều kiện dương \\
10 & 47 & Khoảng cách trong không gian, hay quên bình phương tọa độ \\
\bottomrule
\end{tabularx}}

\subsection{Kịch bản cho mỗi câu (khoảng 4 phút)}

1. \textbf{Gọi học sinh nêu chỗ vướng} (30 giây): không đọc lời giải ngay, hỏi trước "em dừng ở bước nào".
2. \textbf{Chỉ ra từ khóa quyết định} (30 giây): câu này nhìn vào đâu để biết phải làm gì.
3. \textbf{Trình bày đường đi} (2 phút): viết trên bảng theo từng bước, không nhảy bước.
4. \textbf{Cho câu tương tự làm tại chỗ} (1 phút): học sinh tự làm, giáo viên đi quanh lớp.

\subsection{Câu dùng để luyện tại chỗ}

\bang{
\begin{tabularx}{\linewidth}{@{\hspace{6pt}}lX@{\hspace{6pt}}}
\toprule
\textbf{Câu đã chữa} & \textbf{Câu luyện tương tự} \\
\midrule
31 & Tìm giá trị lớn nhất của $y = \dfrac{3x}{x^2 + 1}$ \\
26 & Tìm tập xác định của $f(x) = \sqrt{x - 1} + \dfrac{1}{x - 3}$ \\
34 & Giải bất phương trình $\lvert 3x + 1 \rvert < 5$ \\
40 & Giải bất phương trình $\left(\dfrac{1}{3}\right)^x < 9$ \\
19 & Viết phương trình đường thẳng qua $(1, 3)$ và vuông góc với $2x - 5y + 1 = 0$ \\
24 & Cấp số nhân có $a_2 = 8$, $a_5 = 64$, tìm $a_1$ \\
25 & $a_1 = 3$, $a_{n+1} = 3a_n - 4$, tính $a_4$ \\
28 & Tìm giá trị nhỏ nhất của $f(x) = x^3 - 3x + 2$ trên $[0, 2]$ \\
8 & Cho $\sin\alpha - \cos\alpha = \dfrac{1}{3}$, tính $\sin 2\alpha$ \\
36 & Tìm giá trị nhỏ nhất của $x + \dfrac{9}{x}$ với $x > 0$ \\
\bottomrule
\end{tabularx}}

\section{KHỐI 4 — HỆ THỐNG DẠNG BÀI VÀ KỸ THUẬT XỬ LÝ NHANH (16 phút)}

\subsection{Mục tiêu}

Trước kỳ thi cuối kỳ, học sinh cần một bảng tra duy nhất: nhìn từ khóa, biết ngay dạng bài và cách làm. Khối này chốt lại toàn bộ chương trình trong hai bảng.

\subsection{Bảng hệ thống dạng bài trọng tâm toàn chương trình}

\bang{
\begin{tabularx}{\linewidth}{@{\hspace{6pt}}lllX@{\hspace{6pt}}}
\toprule
\textbf{Nhóm} & \textbf{Dạng bài} & \textbf{Từ khóa nhận dạng} & \textbf{Cách làm} \\
\midrule
Lượng giác & Tính giá trị lượng giác & 已知 $\sin$, 求 $\cos$ & Dùng $\sin^2 + \cos^2 = 1$, xét dấu theo góc phần tư \\
Lượng giác & Rút gọn theo $\tan$ & 已知 $\tan\alpha$, 求分式 & Chia cả tử và mẫu cho $\cos$ \\
Lượng giác & Tổng hợp bậc nhất & $a\sin x + b\cos x$ & Giá trị lớn nhất $\sqrt{a^2+b^2}$ \\
Lượng giác & Chu kì và tập giá trị & 最小正周期, 值域 & $T = \dfrac{2\pi}{\lvert \omega \rvert}$ \\
Hình giải tích & Đường thẳng & 斜率, 截距, 倾斜角 & Rút $y$ theo $x$, đọc trực tiếp \\
Hình giải tích & Đường tròn & 圆心, 半径 & Tâm $\left(-\dfrac{D}{2}, -\dfrac{E}{2}\right)$, $R = \sqrt{\dfrac{D^2+E^2}{4} - F}$ \\
Hình giải tích & Conic & 焦点, 渐近线, 准线 & Đọc $a$, $b$, $c$ từ phương trình chuẩn \\
Dãy số & Cấp số cộng và nhân & 等差数列, 等比数列 & $a_n = a_1 + (n-1)d$, $a_n = a_1 q^{n-1}$ \\
Dãy số & Hệ thức truy hồi & $a_{n+1} = f(a_n)$ & Tính lần lượt, không tìm công thức tổng quát \\
Hàm số & Tập xác định & 定义域 & Căn bậc chẵn $\ge 0$, mẫu $\ne 0$, logarit $> 0$ \\
Hàm số & Đơn điệu và cực trị & 单调递增, 最小值 & Đạo hàm xét dấu, hoặc đỉnh parabol \\
Hàm số & Chẵn lẻ & 奇函数, 偶函数 & Thay $-x$, so với $f(x)$ \\
Tập hợp và bất đẳng thức & Phép toán tập hợp & 交集, 并集, 补集 & Giải về khoảng rồi vẽ trục số \\
Tập hợp và bất đẳng thức & Tham số và tập con & $A \subseteq B$, 恒成立 & Xét $\Delta$, nhớ trường hợp hệ số bằng $0$ \\
Mũ và logarit & Phương trình cùng cơ số & $a^{f(x)} = a^{g(x)}$ & Đưa về cùng cơ số, so số mũ \\
Mũ và logarit & Rút gọn logarit & $\log_a b$, 已知 $\log_a c$ & Tách tích thành tổng \\
Vector và số phức & Tích vô hướng & 垂直, 夹角 & $a \cdot b = 0$ khi vuông góc \\
Vector và số phức & Số phức & 共轭, 虚部, 模 & Nhân liên hợp, dùng $z\bar{z} = \lvert z \rvert^2$ \\
Xác suất và thống kê & Cổ điển & 概率, 随机抽取 & Đếm thuận lợi chia không gian mẫu \\
Xác suất và thống kê & Thống kê mô tả & 平均数, 中位数, 方差 & Công thức trung bình và phương sai \\
Hình không gian & Mặt cầu và mặt phẳng & 球心, 半径, 距离 & Đọc tâm, công thức khoảng cách \\
\bottomrule
\end{tabularx}}

\subsection{Bảng kỹ thuật xử lý nhanh theo nhóm nội dung}

\bang{
\begin{tabularx}{\linewidth}{@{\hspace{6pt}}lXc@{\hspace{6pt}}}
\toprule
\textbf{Nhóm} & \textbf{Kỹ thuật trong phòng thi} & \textbf{Thời gian mục tiêu} \\
\midrule
Lượng giác & Nhớ bảng giá trị góc đặc biệt và dấu theo góc phần tư, không vẽ đường tròn & 40 giây \\
Lượng giác & Thấy phân thức chứa $\sin$ và $\cos$ thì chia ngay cho $\cos$ & 35 giây \\
Hình giải tích & Thay đáp án vào đề khi phương trình phức tạp, đặc biệt câu đường thẳng & 50 giây \\
Hình giải tích & Với conic, viết phương trình chuẩn ra nháp rồi đọc $a$, $b$, $c$ & 45 giây \\
Dãy số & Cấp số thì dùng công thức tổng quát, truy hồi thì bấm máy nhẩm từng bước & 40 giây \\
Hàm số & Tập xác định là câu cho điểm, làm trước tiên nếu gặp & 30 giây \\
Tập hợp và bất đẳng thức & Vẽ trục số, đánh dấu ngoặc tròn cho nghiệm của mẫu & 50 giây \\
Mũ và logarit & Đưa hết về cùng cơ số trước khi biến đổi bất cứ thứ gì & 40 giây \\
Vector và số phức & Số phức nhân liên hợp là con đường ngắn nhất, không tách thủ công & 40 giây \\
Xác suất và thống kê & Liệt kê không gian mẫu bằng bảng nếu số phần tử nhỏ hơn $20$ & 50 giây \\
Hình không gian & Khoảng cách điểm đến mặt phẳng chỉ cần thay số vào công thức & 40 giây \\
\bottomrule
\end{tabularx}}

\subsection{Chiến lược làm bài cho kỳ thi cuối kỳ}

\textbf{Phân bổ thời gian.} 48 câu trong 60 phút, trung bình 75 giây mỗi câu. Chia làm ba vòng:

\bang{
\begin{tabularx}{\linewidth}{@{\hspace{6pt}}lcX@{\hspace{6pt}}}
\toprule
\textbf{Vòng} & \textbf{Thời gian} & \textbf{Nhiệm vụ} \\
\midrule
Vòng 1 & 30 phút & Làm mọi câu chắc chắn đúng, bỏ qua câu chưa nhận ra dạng trong 30 giây \\
Vòng 2 & 20 phút & Quay lại câu đã bỏ, xử lý bằng cách thay đáp án hoặc ước lượng \\
Vòng 3 & 10 phút & Tô lại phiếu trả lời, kiểm tra không bỏ trống câu nào \\
\bottomrule
\end{tabularx}}

\textbf{Ba nguyên tắc bắt buộc.}

1. \textbf{Không dừng quá 90 giây ở một câu.} Hết 90 giây thì đánh dấu và đi tiếp. Câu khó không đáng để mất ba câu dễ.
2. \textbf{Không bao giờ bỏ trống.} Không trừ điểm câu sai, nên luôn tô một đáp án, kể cả khi phải đoán.
3. \textbf{Đọc từ khóa trước, đọc số liệu sau.} Xác định dạng bài rồi mới đọc dữ kiện, tránh bị cuốn vào chi tiết.

\textbf{Thứ tự làm bài đề xuất.} Làm trước các câu nhận ra dạng trong 5 giây: tập xác định, số phức, thống kê, mũ logarit, vector. Để lại cuối các câu hình học giải tích và lượng giác biến đổi dài.

\section{CHỐT BUỔI (5 phút cuối, nằm trong khối 4)}

Viết lên góc bảng và để nguyên khi học sinh ra về:

1. \textbf{75 giây mỗi câu, quá 90 giây thì bỏ qua.}
2. \textbf{Không bỏ trống bất cứ câu nào.}
3. \textbf{Nhận dạng trước, tính toán sau.}

Dặn dò: buổi 20 là buổi ôn tập cuối kỳ, học sinh mang theo toàn bộ bảng đáp án các buổi chữa đề để đối chiếu lại các câu từng sai.

\section{GHI CHÚ SAU BUỔI HỌC}

\emph{(Điền sau khi dạy xong, dùng để điều chỉnh buổi ôn tập cuối kỳ.)}

\bang{
\begin{tabularx}{\linewidth}{@{\hspace{6pt}}Xl@{\hspace{6pt}}}
\toprule
\textbf{Nội dung} & \textbf{Ghi nhận} \\
\midrule
Thời gian thực tế từng khối &  \\
Nhóm nội dung học sinh yếu nhất &  \\
Câu hỏi học sinh hỏi nhiều &  \\
Điều chỉnh cho buổi 20 &  \\
\bottomrule
\end{tabularx}}
\vspace{0.8cm}
\begin{center}
{\footnotesize\color{cae_blue} Tài liệu nội bộ, CAE SHANGHAI}
\end{center}

\end{document}
